\begin{pastexam}[label={px:ncku-106-ee-ds-01}]{成大電機所 106 學年度 複選題第 1 題}
（複選；每題 5 分，全對才給分）

Which of the following statement(s) is (are) true?
\begin{enumerate}[label=(\Alph*)]
\item If $g_1(n) \in \Omega(f_1(n))$ and $g_2(n) \in \Omega(f_2(n))$ then $g_1(n)+g_2(n) \in \Omega(\max\{f_1(n)+f_2(n)\})$
\item $5\log n^2 + \log\log n^2$ is $\Theta(\log n^2)$
\item $f(n) \in \Theta(g(n))$ if and only if $f(n) \in O(g(n))$ and $f(n) \in \Omega(g(n))$
\item $f(n) \in O(g(n))$ and $h(n) \in \Omega(g(n))$ imply $f(n) \in O(h(n))$
\end{enumerate}
\end{pastexam}

\begin{pastsoln}
答案：(A)(B)(C)(D)，四個都對。

照資結書第~1.5~節，函數都是漸近非負的（$n$ 夠大以後 $\ge 0$），這是使用漸近符號的前提。

\textbf{(A) 對。} 題目寫 $\max\{f_1(n) + f_2(n)\}$，大括號裡只有一項，max 就是 $f_1(n) + f_2(n)$；原意
可能是 $\max\{f_1(n), f_2(n)\}$，兩種讀法都對。 由定義，有 $c_1, c_2 > 0$ 和 $n_1, n_2$ 使
$g_1(n) \ge c_1 f_1(n)$（$n \ge n_1$）、$g_2(n) \ge c_2 f_2(n)$（$n \ge n_2$）。 令 $c = \min(c_1, c_2)$，
$n_0 = \max(n_1, n_2)$，$n \ge n_0$ 時
\[
g_1(n) + g_2(n) \ge c_1 f_1(n) + c_2 f_2(n) \ge c\,(f_1(n) + f_2(n)),
\]
所以 $g_1 + g_2 \in \Omega(f_1 + f_2)$。 又 $f_1 + f_2 \ge \max\{f_1, f_2\}$（兩者非負），所以也
$\in \Omega(\max\{f_1, f_2\})$。

\textbf{(B) 對。} $\log n^2 = 2\log n$，$\log\log n^2 = \log(2\log n) = \log 2 + \log\log n$。
\[
\frac{5\log n^2 + \log\log n^2}{\log n^2} = 5 + \frac{\log 2 + \log\log n}{2\log n} \to 5 \quad (n \to \infty),
\]
極限是正的常數，所以是 $\Theta(\log n^2)$（$= \Theta(\log n)$）。 直接用定義（以 $2$ 為底）：$n \ge 2$ 時 $\log n^2 \ge 2$，所以
$0 \le \log\log n^2 \le \log n^2$，下界取 $c_1 = 5$、上界取 $c_2 = 6$、$n_0 = 2$。

\textbf{(C) 對。} 這是 $\Theta$ 的基本性質（資結書第~1.5~節）。 $\Theta$ 的定義是有 $c_1, c_2 > 0$、
$n_0$ 使 $c_1 g(n) \le f(n) \le c_2 g(n)$；左半就是 $\Omega$，右半就是 $O$。 反過來，$O$ 和 $\Omega$ 各給
一組常數和起點，兩個 $n_0$ 取大的那個，兩個不等式同時成立，就是 $\Theta$。

\textbf{(D) 對。} $f \le c_1 g$（$n \ge n_1$），$h \ge c_2 g$，即 $g \le h / c_2$（$n \ge n_2$）。 $n \ge \max(n_1, n_2)$ 時
\[
f(n) \le c_1 g(n) \le \frac{c_1}{c_2}\, h(n),
\]
所以 $f \in O(h)$。 直觀地說：$f$ 不比 $g$ 大，$h$ 不比 $g$ 小，所以 $f$ 不比 $h$ 大。

四個敘述都對，答案 (A)(B)(C)(D)。
\end{pastsoln}
